\documentclass[xcolor=dvipsnames,table]{beamer}

\usepackage{latexsym}
\usepackage[utf8]{inputenc}
\usepackage[brazil]{babel}
\usepackage{amssymb}
\usepackage{amsmath}
\usepackage{stmaryrd}
\usepackage{fancybox}
\usepackage{datetime}
\usepackage[T1]{fontenc}
\usepackage{graphicx}
\usepackage{graphics}
\usepackage{url}
\usepackage{algorithmic}
\usepackage{algorithm}
\usepackage{acronym}
\usepackage{array}

\newtheorem{definicao}{Definio}
\newcommand{\tab}{\hspace*{2em}}

\mode<presentation>
{
  \definecolor{colortexto}{RGB}{0,0,0}
 
  \setbeamertemplate{background canvas}[vertical shading][ bottom=white!10,top=white!10]
  \setbeamercolor{normal text}{fg=colortexto} 

  \usetheme{Warsaw}
}

\title{União, interseção e subgrafos} 

\author{
  Esdras Lins Bispo Jr. \\ \url{bispojr@ufg.br}
  } 
 \institute{
  Teoria de Grafos \\Bacharelado em Ciência da Computação}
\date{\textbf{06 de junho de 2017} }

\logo{\includegraphics[width=1cm]{images/ufgJataiLogo.png}}

\begin{document}

	\begin{frame}
		\titlepage
	\end{frame}

	\AtBeginSection{
		\begin{frame}{Sumário}%[allowframebreaks]{Sumário}
    		\tableofcontents[currentsection]
    		%\tableofcontents[currentsection, hideothersubsections]
		\end{frame}
	}

	\begin{frame}{Plano de Aula}
		\tableofcontents
		%\tableofcontents[hideallsubsections]
	\end{frame}
    
   \begin{frame}{Pensamento}
   	\begin{columns}
   		\column{.4\textwidth}  		
   		\begin{center}
   			\includegraphics[height=.6\textheight]{images/trump.jpg}
   		\end{center}
   		\column{.6\textwidth}  		
   		\begin{block}{Frase}
   			\begin{center}
   				{\large I will build a great wall (and nobody builds walls better than me, believe me)... and I’ll build them very inexpensively. I will build a great, great wall on our southern border, and I will make Mexico pay for that wall. Mark my words.}
   			\end{center}
   		\end{block}		  		
   		\begin{block}{Quem?}
   			\begin{center}
   				{\bf Donald Trump (1946 - ???)} \\Atual presidente estadunidense.
   			\end{center}
   		\end{block}
   	\end{columns}
   \end{frame}
    
    \section{Revisão}
	
	\subsection{Grafo Bipartido (Cont.)}
	
	\begin{frame}{Grafo Bipartido}
		\begin{block}{$K_{p,q}$}
			Um $K_{p,q}$ é um grafo bipartido completo com \\$p$ vértices brancos e $q$ pretos.
		\end{block}
		\begin{block}{Estrela}
			\begin{itemize}
				\item Uma {\bf estrela} é um grafo $K_{1,q}$; 
				\item Se $q \geq 2$, o {\bf centro} da estrela é o único vértice que incide em duas ou mais arestas; 
				\item Se $q < 2$, a estrela não tem centro.
			\end{itemize}
		\end{block}
	\end{frame}

	\subsection{Caminhos e Circuitos}
	\begin{frame}[shrink]{Caminhos e Circuitos}
		\begin{block}{Caminho}
			Um grafo $G$ é um {\bf caminho} se $V_G$ admite uma permutação $(v_1, v_2, \ldots , v_n)$ tal que
			\begin{center}
				$E_G = \{ v_i v_{i +1} : 1 \leq i < n \}$
			\end{center} 
			\begin{itemize}
				\item os vértices $v_1$ e $v_n$ são os {\bf extremos} do caminho; 
				\item os demais vértices são {\bf internos};  
				\item diremos que esse caminho {\bf liga} $v_1$ a $v_n$.
			\end{itemize}
		\end{block}  
		\begin{block}{Notação}
			Podemos denotar um caminho pela sequência representada pelos seus vértices: 
			\begin{center}
				$v_1 v_2 \ldots v_n$		
			\end{center}
		\end{block}
	\end{frame}
	
	\begin{frame}[shrink]{Caminhos e Circuitos}
		\begin{block}{Circuito}
			Um grafo $G$ é um {\bf circuito} se $V_G$ tem 3 ou mais elementos e admite uma permutação $(v_1, v_2, \ldots , v_n)$ tal que
			\begin{center}
				$E_G = \{ v_i v_{i +1} : 1 \leq i < n \} \cup \{ v_1 v_n \}$
			\end{center} 
		\end{block} 
		\begin{block}{Notação}
			\begin{itemize}
				\item Podemos denotar um circuito simplesmente por:  
				\begin{center}
					$v_1 v_2 \ldots v_n v_1$		
				\end{center}
				\item O {\bf comprimento} de um caminho ou circuito $G$ é o número $m(G)$; 
				\item Um {\bf triângulo, quadrado, pentágono {e} hexágono} é o mesmo que um circuito de comprimento 3, 4, 5 e 6 respectivamente.
			\end{itemize}
		\end{block}
	\end{frame}
	\section{União e Intersecção de Grafos}
	\begin{frame}{União e Intersecção de Grafos}
		\begin{block}{União}
			A {\bf união} de dois grafos $G$ e $H$ é o grafo $(V_G \cup V_H, E_G \cup E_H)$. \\É natural denotar esse grafo por $G \cup H$.
		\end{block} \pause
		\begin{block}{Intersecção}
			A {\bf intersecção} de dois grafos $G$ e $H$ é o grafo $(V_G \cap V_H, E_G \cap E_H)$. É natural denotar esse grafo por $G \cap H$.
		\end{block} \pause
		\begin{alertblock}{Alguns cuidados...}
			Para evitar grafos sem vértices, só trataremos da \\interação $G \cap H$ se $V_G \cap V_H$ não for vazio.
		\end{alertblock}
	\end{frame}
	
	\begin{frame}{União e Intersecção de Grafos}
		\begin{block}{Grafos disjuntos}
			Dois grafos $G$ e $H$ são {\bf disjuntos} se \\os conjuntos $V_G$ e $V_H$ são disjuntos.
		\end{block} \pause
		\begin{block}{Corolário}
			Se $G$ e $H$ são disjuntos, então $E_G$ e $E_H$ são disjuntos.
		\end{block}
	\end{frame}

	\section{Subgrafos}
	\begin{frame}{Subgrafos}
		\begin{block}{Definição}
			Um {\bf subgrafo} de um grafo $G$ é qualquer grafo $H$ tal que \\$V_H \subseteq V_G$ e $E_H \subseteq E_G$.
		\end{block} \pause
		\begin{block}{Notações e Nomenclaturas}
			\begin{itemize}
				\item É conveniente escrever ``$H \subseteq G$'' para dizer que \\$H$ é subgrafo de $G$; \pause
				\item Um subgrafo $H$ de $G$ é {\bf gerador} ({\it abrangente}, para alguns) \\se $V_H = V_G$; \pause
				\item Um subgrafo $H$ de $G$ é {\bf próprio} \\se $V_H \not= V_G$ ou $E_H \not= E_G$ (notação: $H \subset G$).
			\end{itemize}
		\end{block}
	\end{frame}
	
	\begin{frame}{Subgrafos}
		\begin{block}{Subgrafo induzido - $G[X]$}
			O subgrafo de $G$ {\bf induzido} por um subconjunto $X$ de $V_G$ é \\o grafo $(X, F)$ em que $F$ é o conjunto $E_G \cap X^{(2)}$. \\Esse subgrafo é denotado por $G[X]$.
		\end{block} \pause
		\begin{block}{$G - X$}
			Para qualquer subconjunto $X$ de $V_G$, \\denotaremos por $G - X$ o subgrafo $G[V_G \setminus X]$.
		\end{block} \pause
		\begin{block}{$G - v$}
			Uma abreviação para $G - \{ v \}$.
		\end{block}
	\end{frame}
	
	\begin{frame}{Subgrafos}
		\begin{block}{$G - a$}
			Uma abreviação para o grafo $(V_G, E_G \setminus \{ a \})$.
		\end{block} \pause
		\begin{block}{$G - A$}
			Se $A$ é um subconjunto de $E_G$, então $G - A$ é uma abreviação para o grafo $(V_G, E_G \setminus A)$.
		\end{block} \pause
		\begin{block}{Corolário}
			$G - A$ é um grafo gerador de $G$.
		\end{block}
	\end{frame}
	
	\begin{frame}
		\titlepage
	\end{frame}
	
\end{document}